\chapter*{Conclusion}
\addcontentsline{toc}{chapter}{Conclusion}

So, can an operator just ask a question in plain language and get back an answer
that is actually correct, actually current, and honest about what it does not
know? After everything covered in this report, the answer is yes --- but with
conditions, and honestly the conditions turned out to be the most interesting
part of the whole thing.

What we ended up with is a working system, not a demo of one. Two databases
modelled on TM~Forum's abstractions holding around 50{,}000 subscribers, a
simulator that keeps them moving --- raising faults, degrading the cells those
faults sit on, pushing the damage down to the customers that cell serves, then
healing it all again --- and an agent that answers questions against that moving
target, grounds every query in the real schema instead of guessing, checks its
own output before handing it over, and refuses to touch anything operational
without a human saying yes first.

Measured on a golden set of sixteen questions, it gets every factual question
right, refuses every destructive instruction and every prompt injection, and
does not fall over on questions that have no answer at all. It fails at exactly
one thing, and it fails at it reliably: admitting that it does not have
something. Ask it for a metric the schema simply does not hold, and instead of
saying so, it grabs the closest real metric it can find and reports that
instead. Every number in that answer is real, which is exactly why every
verification layer in the system waves it straight through.

Three things came out of this project that are worth carrying beyond it.

\textbf{The model was never the bottleneck.} Given the same questions, a much
larger model made the same kinds of mistakes as the smaller one: joining through
the wrong table, claiming things it had not measured, quoting offer terms it
never looked up. That is not a reasoning problem. It is the model not knowing
some fact about the schema, or the code around it letting an unchecked claim
slip through. Which is why most of this report ended up being about the
machinery wrapped around the model rather than the model itself.

\textbf{Correctness has to be forced, not asked for.} The system prompt already
told the agent to claim only what it had queried, and it still went and said
nothing else was affected while a much bigger outage was running somewhere else.
What actually worked was a check that runs afterwards and does not care what the
prompt said: pull out the claims, ask whether a query backs each one, rewrite
the answer from the real rows when it does not. Telling the model something is a
request. Checking it afterwards is a guarantee --- but only for the specific
kind of error you built the check for, which is the whole lesson.

\textbf{A number measured on one dataset tells you nothing about another one.}
The churn model reported 0.88 AUC-ROC, and it was pointed at an environment
whose feature distributions differ from its training data by up to a factor of
three hundred. Both of those were true at the same time, for months, and nobody
noticed because nobody checked. We swapped it for a scorecard built from the
operator's own columns --- which is not validated either, and this report says
so rather than quietly hoping nobody asks.

Underneath all three of those sits something this project kept running into:
systems that report success while doing absolutely nothing. A KPI updater that
had never written a single row. A table that stopped moving because its backfill
window was anchored to the data instead of the clock. An evaluation harness that
had never once run, because nothing ever loaded its credentials. Not one of them
threw an error. They all just produced confident, stale answers, which is much
harder to spot than something that crashes --- and learning to spot that was
probably worth more to me than any feature I added.

As for the objectives set out at the start, they were met, with one exception
worth stating plainly: the system got built, made live, grounded, and kept on a
leash so it proposes instead of acting --- but the target of 80\% correctness on
a golden set was not reached, and the evaluation that established that is narrow
enough that the number should be read as an indication rather than a
measurement. Chapter~\ref{ch:limits} goes through what that leaves undone, in
the order it would actually be worth doing.

In the end this project taught me less about telecom analytics than about
building anything on top of a language model. Getting a model to produce a
plausible answer was never the hard part. The hard part was building enough
machinery around it to know whether that answer was true --- and then finding
out, by finally measuring it, that the machinery catches less than it looks like
it does.
